\chapter{Diskusjon – funn og kobling til forskningsspørsmål}

I dette kapittelet diskuteres de empiriske funnene fra brukertestingen, og de knyttes til forskningsspørsmålene (RQ1–RQ4) og til problemstillingen. Diskusjonen hviler på to datakilder: to runder med anonyme spørreundersøkelser på en 1–7 Likert-skala (kvantitativt), gjennomført før og etter designendringene, og fem semistrukturerte dybdeintervjuer (kvalitativt). Tenke-høyt-opptakene fra de samme sesjonene analyseres i et senere ledd og vil supplere de kvalitative funnene. Kapitlet er bygd opp slik at det først redegjøres for (1) hvilke resultater de metodiske valgene ga, deretter (2) hvordan resultatene tolkes i lys av teorikapitlet, og til slutt (3) hvordan hvert enkelt forskningsspørsmål besvares med utgangspunkt i problemstillingen \cite{weichbroth2024,ericsson1993,braun2006}.
\newline 

\section{Hva de metodiske valgene ga}
Denne studien er utformet som en evalueringsstudie av brukervennlighet, og det ble gjennomført tester basert på scenarier, tenke-høyt-teknikken og semistrukturerte intervjuer \cite{weichbroth2024,ericsson1993}. Det ble tatt hensyn til valg av metode for å avdekke både:
(a) observerbare atferder og feil i registreringsprosessen, og
(b) brukerens mentale modeller og forventninger knyttet til lagring, innsending og tilbakestilling.
\newline
Studien ble gjennomført som en evaluativ brukbarhetstest i to faser, ved å teste to av anonyme spørreundersøkelser på en 1–7 Likert-skala. Første testing (25.–26. februar 2026, 31 testere) dokumenterte utgangspunktet før vi implementerte designendringer basert på de mest kritiske tilbakemeldingene. For å kunne måle effekten av de tilbakemeldinger som er samlt fra testere, ble en identisk undersøkelse fullført mellom 27. februar og 10. mars 2026 (med 21 person). De samlede resultatene viste en klar forbedring: Samlet appscore for Q2–Q10 steg fra 77,7\% til 86,2\% (+8,5 prosentpoeng), og fullføringsraten forble på 100\% gjennom hele studien.

De tre metodene utfyller hverandre og er rettet mot hvert sitt aspekt av registreringsflyten:
\begin{itemize}
  \item \textbf{Scenario-basert brukertesting} er innrettet mot å fange opp punktene der brukerne stopper opp, avviker fra prosessen eller avbryter den, eller utfører handlinger som kan skape inkonsistens i dataene (for eksempel gjentatt innsending, feil rekkefølge eller å hoppe over trinn). Hensikten er å fremkalle de kritiske tilstandsendringene i systemet: lagring, navigasjon, innsending og tilbakestilling \cite{weichbroth2024}.
  \item \textbf{Tenke-høyt-teknikken} gir direkte innsikt i hva brukerne tror at systemet «husker», hva de forventer blir lagret, og hvordan de tolker systemtilstanden («lagret», «sendt», «slettet»). Den utgjør grunnlaget for å identifisere manglende samsvar i den mentale modellen (\textit{mental model mismatch}). Disse opptakene analyseres i et senere ledd \cite{ericsson1993,norman2013}.
  \item \textbf{De semistrukturerte intervjuene} gir etterfølgende begrunnelse og klargjøring: hvorfor brukeren utførte en bestemt handling, hvilke signaler i brukergrensesnittet som ble oversett, og hvilke forbedringer brukeren selv mener vil gjøre arbeidsflyten mer forutsigbar \cite{weichbroth2024}.
\end{itemize}

I den tematiske analysen organiseres funnene ved hjelp av koder som samsvarer med definisjonene i metodekapitlet: \textit{uklar systemtilstand}, \textit{persistensproblemer}, \textit{resetproblemer}, \textit{mental model mismatch} og \textit{datainkonsistens} \cite{braun2006}.
\newline




\section{Teoretisk forankring for diskusjonen}

Diskusjonen vil bygges på tre sentrale teoretiske konsepter som tidligere er forklart i oppgaven::
\begin{enumerate}
  \item \textbf{Systemtilstand og status-synlighet}: Når systemets interne tilstand (utkast, lagring, innsending, tilbakestilling) ikke er klar eller konsistent, øker risikoen for feil og misforståelser \cite{nielsen1994}.
  \item \textbf{Mentale modeller}: Brukere handler basert på antakelser om hvordan systemet fungerer.
Manglende samsvar mellom den mentale modellen og systemets faktiske funksjonalitet fører til friksjon, feil og redusert tillit \cite{norman2013}.
  \item \textbf{Datakvalitet/datakonsistens}: Kvalitet forstås her som noe som går utover nøyaktighet; fullstendighet, unikhet og korrekt kobling (bilde – metadata) er nødvendige for fremtidig bruk \cite{wang1996}.
\end{enumerate}

Forskningsmodellen i oppgaven presenterer en tydelig kausal kjede som brukes i drøftingen for å forklare funnene: \emph{systemtilstand} $\rightarrow$ \emph{status-synlighet} $\rightarrow$ \emph{mental modell} $\rightarrow$ \emph{feilrate} $\rightarrow$ \emph{datakonsistens}. Denne modellen brukes som et rammeverk for å knytte brukeropplevelsesproblemer til konkrete former for datainkonsistens i registreringsprosessen \cite{nielsen1994,norman2013,wang1996}.

\section{Diskusjon og ``funn''-typer koblet til forskningsspørsmål}

For å belyse problemstillingen er det formulert fire forskningsspørsmål som dekker ulike aspekter av brukeropplevelse, systemtilstand og datakvalitet i registreringsflyten. Spørsmålene er utformet med utgangspunkt i det teoretiske rammeverket og bygger på antakelsen om at svakheter i brukergrensesnitt og teknisk implementasjon kan føre til feil og inkonsistente data. I dette kapittelet presenteres hvert forskningsspørsmål med tilhørende datakilder, forventede funntyper, teoretisk forankring og hvordan det bidrar til å besvare problemstillingen.


\subsection{RQ1:- Hvordan opplever ansatte med ulik teknisk kompetanse brukervennligheten i registreringsflyten (bilde → kategori → metadata → innsending), spesielt med tanke på forutsigbarhet, feil og fullføring?
}

\textbf{Funnkilder:} Observasjoner under gjennomføring av scenariobaserte oppgaver, tenke-høyt-data og intervjusvar om flyt, kontroll og forutsigbarhet, samt kvantitative resultater fra to brukerundersøkelser gjennomført før og etter designendringer.

\textbf{Forventede funntyper.}
\begin{itemize}
  \item  Punkter i arbeidsflyten der brukerne føler usikkerhet ("Hvor er jeg i prosessen?").
  \item Situasjoner der brukerne søker bekreftelse på hva som er lagret/sendt.
  \item Ulike mønstre mellom tekniske og ikke‑tekniske deltakere (for eksempel at tekniske deltakere tester systemets grenser, mens ikke‑tekniske deltakere følger grensesnittet bokstavelig).
\end{itemize}

\textbf{Funn – empiriske resultater: } Dataene som ble samlet inn, bekrefter at designendringene som ble gjort i applikasjonen, har bidratt til å gjøre arbeidsflyten i registreringsprosessen tydeligere. Dette bidro til at vurderingen av applikasjonen (Q2–Q10) økte med 8,5 prosentpoeng (fra 77,7\% til 86,2\%). Samtidig forble fullføringsgraden stabil på 100\% i begge rundene der testingen ble gjennomført. På bakgrunn av dette kan vi konkludere med at den grunnleggende registreringsstrukturen ikke utgjorde en hindring, men at variasjonen lå i kvaliteten på brukergrensesnittet.

Resultatene i første runde viste en god forståelse av trinnrekkefølgen i prosessen (bilde → kategori → metadata → innsending), med 86,6\% på spørsmål 2 (Q2). Samtidig ligger det mest fremtredende resultatet for besvarelsen av det første forskningsspørsmålet (RQ1) i det markante fallet i resultatet for spørsmål 8 (Q8), som bare nådde 39,2\%. Dette tallet viser en svakhet i tydeligheten til ikonene som ble brukt, samt i selve registreringsprosessen. Denne svakheten anses som alvorlig i mobilapplikasjoner som brukes av brukere under krevende arbeidsforhold som stiller krav til nøyaktighet og hurtighet i utførelsen \cite{muller2024, wenz2023}; manglende evne til å skille mellom knappene strider mot prinsippene for brukervennlighet som ble fastsatt av Wichbroth\cite{weichbroth2024}, som i hovedsak bygger på effektivitet og reduksjon av kognitiv belastning.

De kvantitative dataene, sammen med de kvalitative kommentarene som ble samlet inn i første runde, bekrefter dette, ettersom 26 utsagn (omtrent halvparten) inneholdt klager på manglende tekstlige etiketter for ikonene, slik som: «Ikonene har ikke tekst, så det bør være tekst under dem», «sendeknappen og de andre ikonene (ta bilde eller laste opp) har ingen tekstlig forklaring», og «det var ikoner som irriterte meg fordi det ikke fantes tekst under dem». Dette er en indikator på i hvilken grad brukere med ikke-teknisk bakgrunn er avhengige av eksplisitte veiledninger i brukergrensesnittet. Etter endringen som ble gjort i applikasjonen ved å legge til tekst, steg Q8 til 91,2\% (en økning på 52,0 prosentpoeng). For å sikre påliteligheten av dette resultatet viste kontrollvariabelen Q1 (teknisk erfaring) stabilitet mellom de to rundene (61,3\% mot 59,2\%), noe som isolerer effekten av brukernes ferdigheter og bekrefter effekten av designendringen.

\newline 
\textbf{Funn fra dybdeintervjuene: }
Resultatene fra de kvalitative intervjuene som ble gjennomført med flere deltakere, gjenspeiler en styrking av de positive indikatorene som ble registrert i de kvantitative spørreundersøkelsene. De fem deltakerne var enige om hvor smidig og tydelig registreringsprosessen var. Derfor beskrev både  «Jeg synes det var veldig enkelt og greit» og «Ingenting var vanskelig. Jeg synes det var veldig, veldig enkelt og greit» og «Erik» fra Stena Recycling AS sa «Jeg synes det var veldig, veldig oversiktlig» , og bekreftet dermed at den var fri for nevneverdige hindringer. Dette ble støttet av «Benjamin», som beskrev den som en prosess som «i realiteten er enkel», mens «NG Roger AS» oppsummerte applikasjonen med at «den er brukervennlig». Disse utsagnene gjenspeiler den merkbare forbedringen i de generelle holdningene som spørreundersøkelsene registrerte mellom første og andre runde.

Til tross for dette ga intervjuene mulighet til å avdekke presise detaljer som spørreundersøkelsene ikke klarte å berøre. Selv om det generelle inntrykket var enkelhet, var brukerne ikke sikre når det gjaldt de fullførte oppgavene. «Erik fra Stena Recycling AS» pekte på en uklarhet knyttet til hvordan applikasjonen ser ut til å ha avsluttet registreringsprosessen for saken før den faktisk er sendt, mens «Roger Aas fra NG» hadde vansker med å finne riktig kategori på grunn av fraværet av alternativet «restavfall». Disse observasjonene peker på en viktig metodisk konklusjon, nemlig at «fullføringskvalitet» (Quality of Completion) strukturelt skiller seg fra «fullføringsrate» (Completion Rate). Brukeren kan lykkes med å fullføre alle trinnene prosedyremessig, men likevel ende opp med å registrere feil data.

Intervjuene som ble gjennomført, avdekket også en tydelig variasjon i brukernes behov basert på deres arbeidsroller og arbeidskontekster. «Franser fra Stena Recycling AS» uttrykte stor tilfredshet med hvor enkelt applikasjonen var å bruke, mens «Roger Aas fra NG» presenterte et kritisk perspektiv på brukervennligheten, ved at systemet ikke når opp til de komplekse daglige oppgavene de arbeider med. Både «Erik fra Stena Recycling AS» og «Werner Reiersen fra Stena Recycling AS» gikk dypere inn i krav til avanserte arbeidsflyter, som håndtering av farlig avfall, integrasjon med systemer for virksomhetsressursplanlegging (ERP), programmeringsgrensesnitt (API) og arbeid i miljøer uten behov for internettforbindelse, mens «Benjamin» fokuserte på administrative og analytiske aspekter som søk, filtrering og statistikk.

Denne variasjonen samsvarer med de teoretiske forventningene om ulike bruksmønstre, men den viser at den grunnleggende variabelen i denne variasjonen ikke er brukerens «tekniske kompetanse» slik den ble målt i spørsmål én, men «arbeidskonteksten» og «arten av de funksjonelle arbeidsoppgavene». Dette er et presist kvalitativt skille som de kvantitative spørreundersøkelsene ikke klarte å registrere selvstendig.


\newline 
\textbf{Teorikobling: } Når disse resultatene analyseres i lys av «tydelighet og konsistens i systemets tilstand», blir det klart at utformingen av grensesnittet i en flertrinns arbeidsflyt spiller en avgjørende rolle i hvordan brukeren veiledes \cite{nielsen1994}. Den kvalitative forbedringen i resultatene for spørsmål åtte (Q8) gir en empirisk bekreftelse på Nielsens regel som foretrekker gjenkjenning av symbolet fremfor å måtte huske det \cite{nielsen1994}; i stedet for å la brukeren måtte anstrenge seg for å huske funksjonen til et ikon, bør den presenteres tydelig i grensesnittet.

Fra en annen vinkel kan den lille nedgangen som oppstod i spørsmål to (-3,6 prosentpoeng) forklares som et naturlig fenomen knyttet til de «mentale modellene» som Norman beskrev \cite{norman2013}; da han påpekte at brukere som er vant til den gamle flyten, har en tendens til å oppleve midlertidig forvirring når de møter en ny oppdatering, selv om det nye designet teknisk sett er bedre.
Ved å se på Normans rammeverk \cite{norman2013} bredere i intervjudataene, kan vi tydelig se at applikasjonens grunnleggende flyt faktisk samsvarte med brukernes mentale modeller, men at det oppstod svakheter i «separasjonen» mellom kategoriene og metadatafeltene på den ene siden og den kontekstuelle modellen for faktiske praksiser innen avfallshåndtering på den andre. Dette samsvarer også med observasjonen til Müller mfl. \cite{muller2024} om at brukervennlighet i datainnsamlingsapplikasjoner ikke bare handler om visuell tydelighet, men i hovedsak avhenger av i hvilken grad systemet samsvarer med arbeidssituasjoner i den virkelige verden og deres profesjonelle kontekster.

\newline
\textbf{Bidrag til problemstilling.} Resultatene fra det første spørsmålet bekrefter empirisk at det å redusere sannsynligheten for feil under registrering henger nært sammen med brukerens evne til å «forutsi» hva som vil skje i grensesnittet. Siden uklare symboler var den viktigste kilden til forvirring i første runde, bidro endringen av disse gjennom tilføring av forklarende tekster til å redusere «friksjon» og heve kvaliteten på utførelsen, noe som la grunnlaget for en konsistent brukeropplevelse som faktisk er brukbar i praksis.
Intervjuene også viser imidlertid at kvaliteten på prestasjonen ikke bare avhenger av symbolenes tydelighet, men også av tilpasning til konteksten, et tema som diskuteres mer detaljert i den tredje forskningshypotesen.

\subsection{RQ2: Hvordan påvirker utformingen av brukergrensesnittet (f.eks. tydelige knapper, tilbakemeldinger og konsistens) brukerens forståelse av systemets tilstand – inkludert hva som er lagret, sendt og nullstilt?}

\textbf{Funnkilder.} Disse kvantitative resultatene er sentrert rundt to anonyme brukertester som ble gjennomført før og etter designendringene som ble gjort i applikasjonen (med deltakelse fra henholdsvis 31 og 21 brukere), i tillegg til fritekstkommentarer fra begge undersøkelsene. Disse resultatene vil i en senere analyse bli supplert med tenke-høyt-data fra fem dybdeintervjuer med Franser fra Stena Recycling AS, Erik fra Stena Recycling AS, Benjamin, Werner Reiersen fra Stena Recycling AS og Roger Aas fra NG (Norsk Gjenvinning).

\textbf{Funn – teoretisk bakgrunn: } Forventningen i denne studien var at det ville bli avdekket tre typer misforståelser knyttet til systemets tilstand. Den første var forventningen om datakontinuitet: Brukeren forventer at applikasjonen oppbevarer metadataene selv om det navigeres mellom sidene. Den andre var forventningen om tilbakestilling: Etter at bilder og informasjon er sendt, forventer brukeren at siden blir tilbakestilt og at informasjonen som er sendt, blir slettet. Den tredje var forventningen om fullføring av prosessen: Brukeren forventer et tydelig sluttpunkt som viser at registreringen er fullført, og at applikasjonen er klar for neste oppgave.


\textbf{Forventede funntyper.}
\begin{itemize}
  \item \textbf{Persistensforventning}:  Her vil brukeren forvente at metadata forblir lagret etter navigasjon mellom sider.

  \item \textbf{Resetforventning}: Her vil brukeren forvente en tom tilstand etter at bilder og informasjon er sendt inn.
  \item \textbf{Transaksjonsforventning}: Brukeren forventer et tydelig sluttpunkt («for eksempel ved registrering av en ny konto bør man automatisk gå til innloggingssiden»).
\end{itemize}

\textbf{Funn – empiriske resultater: } De tre dimensjonene som viste sterk forbedring i testen, demonstrerte hver på sin måte systemets evne til å formidle systemets tilstand tydelig til brukeren, og dette bekrefter empirisk de tre typene misforståelser som på forhånd var forventet.

\newline 
Q7 (tydelig fullført etter innsending) økte fra 4,74 (67,7\%) til 6,33 (90,5\%), noe som tilsvarer en økning på 22,7 prosentpoeng. Andelen på 67,7\% i første runde viser at mer enn en tredjedel av deltakerne i testen ikke var sikre på om registreringsprosessen faktisk var fullført eller ikke, og dette utgjør en direkte empirisk bekreftelse på de teoretiske forventningene. Når brukeren ikke vet om registreringen er sendt eller ikke, vil vedkommende kunne sende den flere ganger, noe som fører til dupliserte data, eller forlate systemet uten å fullføre innsendingen, noe som vil føre til generering av ufullstendige data \cite{wang1996}. Etter at bekreftelsen på innsending ble forbedret, steg andelen til 90,5\%, noe som indikerer at tiltakene som ble gjennomført, behandlet denne feilkilden direkte og effektivt.

\newline 
Q3 (bildeopplasting uten tap av kvalitet) økte fra 5,03 (71,9\%) til 6,24 (89,1\%), altså med en økning på 17,2 prosentpoeng. Sju av 26 fritekstkommentarer i første runde nevnte også at tidligere sendte bilder fortsatt ble vist – et vanlig teknisk problem som oppstår når det ikke finnes en mekanisme som tilbakestiller systemet etter innsending. Dette samsvarer fullt ut med den antatte forventningen om tilbakestilling. I Normans terminologi \cite{norman2013} utgjør dette et brudd på systemets konseptuelle modell: Her oppsto det faktisk forvirring om hva som faktisk tilhørte den pågående registreringen, fordi grensesnittet viste en tilstand som ikke gjenspeilte systemets faktiske tilstand.

\newline 
Q9 (trygghet via meldinger) økte også fra 5,74 (82,0\%) til 6,24 (89,1\%), altså med en merkbar forbedring på 7,1 prosentpoeng. Dette problemet gjenspeiler i hvilken grad systemet er i stand til å opprettholde kontinuerlig kommunikasjon med brukeren gjennom statusmeldinger og tilbakemeldinger, og denne forbedringen viser at endringene som ble innført i applikasjonens verifikasjonsmekanismer, bidro til å øke tilliten til at dataene behandles korrekt \cite{wenz2023}.

\newline
\textbf{Funn fra dybdeintervjuene: }
Intervjuene viser at endringene som ble innført i bekreftelsen av innsending, oppnådde det ønskede formålet og gir en presis beskrivelse av brukernes nåværende opplevelse. Franser fra Stena Recycling AS beskrev det som skjer etter innsending slik: «Det ble sendt. Greit, da er den klar for neste innsending», og spurte deretter: «Hva skjer etter innsending?» Erik fra Stena Recycling AS sa: «Jeg visste at det ble sendt da popup-vinduet dukket opp». Benjamin bemerket: «Det dukket opp en liten rød indikator formet som en bjelle nederst til venstre». Werner Reiersen fra Stena Recycling AS oppsummerte det slik: «Ja, dette er veldig tydelig». Brukerne forstår nå at innsendingen er gjennomført, og dette er et empirisk bevis på at forventningene deres knyttet til transaksjoner på visuelt nivå er oppfylt.

Likevel avdekket intervjuene en forskjell som spørreundersøkelsene ikke dekket: Brukerne vet at noe er sendt, men de vet ikke hvor det er sendt, eller hvem som kommer til å se det. Her handler det om forskjellen mellom lokal tilbakemelding og tillit til systemet. Franser fra Stena Recycling AS sa ærlig: «Nei, når jeg først har sendt dataene, vet jeg ikke hvem som kan se dem», og la til: «Det handler om sikkerhet, ikke sant?». Erik fra Stena Recycling AS var enda tydeligere: «Det stemmer, jeg har ingen anelse om hva som skjedde med dataene jeg sendte». Dette kan muligens forklare den svake nedgangen på 3,3\% i Q10-indikatoren (opplevd sikkerhet) etter endringene: Endringene i brukergrensesnittet løste ikke problemet med datatransparens, men kan snarere ha økt usikkerheten ved å flytte fokuset fra «fungerer dette systemet i det hele tatt?» til «hva skjer med dataene etter innsending?».

Endringene forbedret tydeligheten i innsendingen, men reiste samtidig et nytt spørsmål som ikke tidligere var stilt: Hva skjer med dataene etterpå? Med andre ord var bekreftelsen på innsending vellykket, men tillit til systemet er et langt dypere problem enn dette.


\textbf{Teorikobling.} Resultatene som ble samlet inn gjennom disse testene, kan forklares som et tilfelle av manglende samsvar i den mentale modellen, gjennom systemets konseptuelle rammeverk som ikke fremstår tydelig gjennom signalene i brukergrensesnittet og tilbakemeldingene \cite{nielsen1994, norman2013}. Forbedringen som ble observert i spørsmål syv, utgjør også et empirisk bevis på Nielsens heuristiske regel om synlighet av systemstatus \cite{nielsen1994}, som understreker nødvendigheten av at systemet alltid gir brukeren tydelig informasjon om hva som skjer. Utviklingen i spørsmål tre bekrefter dessuten at den tekniske implementeringen av tilstandshåndtering og den visuelle representasjonen av denne i grensesnittet er to sammenknyttede elementer som ikke kan skilles fra hverandre.

\textbf{Bidrag til problemstilling.} RQ2 - Det viser \emph{hvordan} et mulig misforhold i brukerens forståelse av lagring og gjenoppretting kan føre til tap av metadata, duplisering av data eller kommunikasjonsfeil utilsiktet. De empiriske resultatene viste at endringer i brukergrensesnittet, som gjør systemets tilstand tydelig og transparent, kan spille en viktig rolle i å redusere feilkilder og dermed styrke datakonsistens og datakvalitet. Intervjuene viser også at den overordnede løsningen også bør fokusere på hvordan åpenhet bidrar til behandlingen av etterfølgende data.

\subsection{RQ3: - Hvilke typer brukerfeil og datainkonsistenser oppstår i registreringsflyten, og hvordan henger disse sammen med systemets håndtering av tilstand og brukerens mentale modell?}

\textbf{Funnkilder.} Fritekstkommentarer fra begge anonyme brukerundersøkelser (26 kommentarer i første undersøkelse, 17 i andre undersøkelse) og kvantitative endringer per spørsmål. Think-aloud-data fra fem dybdeintervjuer vil supplere med observerte feilhandlinger og årsaksforklaringer i en etterfølgende analyse.

\textbf{Funn – teoretisk bakgrunn: } Med utgangspunkt i gjennomgangen av det teoretiske rammeverket og anvendelsen av studien ble det avdekket en forventning om fire mulige typer feil som påvirker registreringsprosessen.
\begin{itemize}
    \item Tap av persistens i metadataene (Persistence Error): Denne tilstanden kommer til uttrykk ved at de innskrevne dataene forsvinner ved bytte mellom skjermer eller ved navigasjon mellom dem, noe som fører til at innsatsen brukeren har lagt ned i å registrere denne informasjonen, går tapt.

    \item Dataoverlapping eller «lekkasje» (Reset Error): Denne tilstanden viser til at tidligere data dukker opp i en ny registrering som følge av en feil i systemet.

    \item Manglende samsvar mellom bildet og metadataene (koblingsfeil): Dette innebærer at bildet ikke samsvarer med den aktive registreringen.

    \item Uklarhet i den nåværende systemtilstanden (Feedback Ambiguity): Denne tilstanden er knyttet til dobbel innsending eller unødvendige kontrollhandlinger.
\end{itemize}


\textbf{Forventede funntyper (i tråd med operasjonaliseringen).}
\begin{itemize}
  \item \textbf{Metadata-tap} (persistensfeil): data forsvinner ved skjermbytte/navigasjon.
  \item \textbf{Ghost-data/lekkasje} (resetfeil): tidligere data vises i ny registrering.
  \item \textbf{Feilkobling bilde--metadata} (koblingsfeil): feil sammenstilling i aktiv registrering.
  \item \textbf{Statusuklarhet} (``sendt'' er uklart): kan gi dobbel innsending eller unødvendige kontrollhandlinger.
\end{itemize}

\textbf{Funn – empiriske resultater: } De tekstlige kommentarene (fritekstkommentarene) som ble samlet inn gjennom de to rundene av spørreundersøkelsen, ga empirisk støtte til alle de fire feiltypene som var forventet i planleggingsfasen, noe som ga dem et konkret innhold.
\begin{enumerate}
    \item Uklarhet i systemets tilstand (State Ambiguity): Uklarheten i denne tilstanden ble dokumentert gjennom de innledende resultatene i spørsmål sju (Q7), som var på 67,7 \%, og deltakernes kommentarer gjenspeilte også en manglende evne til å fastslå navigasjonsforløpet og fraværet av visuelle signaler, slik som: «Jeg visste ikke hvor sendeknappen var, så jeg måtte stoppe innsendingen» og «Knappene er uklare, og det var vanskelig å vite hva forrige steg var». Dette utgjør en nær sammenheng mellom denne typen feil og det Norman \cite{norman2013} beskrev i lys av begrepet «kunnskap i verden»: Når informasjon om handlingen ikke er tilgjengelig i det visuelle grensesnittet, tyr brukerne til gjetning eller avbrytelse.

    \item Datalekkasje og dataoverlapping (Data Leakage \& Phantom Data): Studien avdekket også i sju kommentarer fra første runde tilstedeværelsen av fantomdata og datalekkasje gjennom utsagn som: «Bildene slettes ikke etter at de er sendt og vises igjen», «De sendte bildene vises igjen når et nytt bilde tas», og «Bildene gjentas når de sendes, selv om det egentlig skulle tas et nytt bilde». I lys av Wang og Strongs rammeverk for datakvalitet \cite{wang1996} kan denne feilen påvirke dimensjonen korrekt tilknytning negativt: Når det oppstår et teknisk avvik der det viste bildet tilhører en prosess som allerede er fullført, men tilsynelatende er koblet til nye metadata, undergraves registreringens pålitelighet.

    \item Datatap som følge av navigasjon (Navigation-Induced Data Loss): Denne tilstanden viste til sletting av kvalitative data ved bakovernavigasjon (å gå fram og tilbake mellom sider), slik som i følgende kommentar: «Hvis jeg skriver inn ordrenummer og kundenummer og navigerer fram og tilbake, mister jeg informasjonen». Dette resultatet samsvarer med konklusjonene til «Müller og andre» \cite{muller2024}, som betrakter dette som en av de mest fremtredende årsakene til datamangel i systemer for kollektiv informasjonsinnsamling (Crowdsourcing), og det er knyttet til implementeringen av lagring som er beskrevet i avsnitt 3.4.

    \item Kategoristøy (Categorization Noise) – det fremvoksende mønsteret: Her fremtrer en fjerde, ny og uventet type feil i de senere dataene, i motsetning til de tidligere mønstrene: støy i kategorivelgeren, der tastaturet som vises under skriving, førte til at teksten som ble skrevet inn, ble skjult, mens de interaktive bevegelsene (animations) i elementene for kategorivalg skapte forvirring hos brukerne i stedet for å bekrefte valgene deres. Blant kommentarene var: «Endring av kategori i bildet var skjult bak tastaturet, så jeg kunne ikke lese det jeg skrev. Kategoriene hopper når de velges, noe som gjør det vanskelig å vite om jeg har utført riktig handling». Denne feilen er trolig ansvarlig for nedgangen i resultatene for spørsmål fire (Q4) på 10,8 prosentpoeng. Dette resultatet belyser at iterativ redesign kan introdusere nye hindringer for brukervennlighet, selv om målet kun er forbedring. Når det gjelder datakvalitet, innebærer dette en økt risiko for feilklassifisering, noe som direkte reduserer nøyaktigheten i treningsdataene \cite{geiger2021}.
\end{enumerate}


\newline
\textbf{Funn fra dybdeintervjuene: }
Intervjuene begrenset seg ikke bare til å bekrefte feiltypene som spørreundersøkelsene hadde registrert, men avdekket også et mer alvorlig mønster som består i «kontekstuell feilklassifisering»; der brukeren blir nødt til å velge en kategori bare fordi den finnes i listen, selv om den ikke samsvarer med den faktiske virkeligheten for avfallet.
Alle deltakerne var enige om dette klassifiseringsproblemet. Benjamin opplyste dermed at: «Kategorien kan variere litt, og kanskje finnes det svært brede kategorier sammenlignet med det du har». På sin side var Werner Reiersen fra Stena Recycling AS tydeligere da han sa: «For meg var alle kategoriene litt uklare», og forklarte dette med et konkret eksempel: «Som med batterier, vil batterier og farlig avfall bli behandlet på samme måte». Roger Aas fra fra NG(Norsk Gjenvinning) påpekte også en grunnleggende mangel som består i fraværet av kategorien «restavfall», og bekreftet at den: «faktisk er en av de største avfallskategoriene», noe som betyr at applikasjonen mangler en av de viktigste avfallskategoriene som håndteres i deres daglige arbeidsmiljø.
Disse svært viktige observasjonene utgjør en direkte trussel mot datakvaliteten i henhold til rammeverket til Wang og Strong \cite{wang1996}, og viser spesielt til dimensjonene «egnethet» og «datamengde»; kategoriene er enten for brede og oppfyller ikke presis klassifisering, eller de mangler helt. Ved å anvende perspektivet til Geiger. \cite{geiger2021} om «dårlige input» blir det klart at kunstig intelligens-modellen vil bli trent på feil data som er lagt inn systematisk, ikke fordi brukeren har lagt inn data feil, men fordi systemet mangler evnen til å tilby de riktige alternativene.
Intervjuene avdekket også at dette gapet også omfatter metadata. Benjamin snakket om sin usikkerhet rundt begrepene ved å si: «Jeg er litt usikker på detaljene i bestillingen og kunde-ID og så videre, fordi jeg vanligvis ikke gjør dette», og foreslo at: «Vi trenger kanskje opplæring». Roger Aas spurte kritisk: «Bestillingsdetaljer. Hva mener du at jeg skal bruke her? Bestillingsdetaljer. Veinummer, for eksempel». Franser fra Stena Recycling AS pekte også på et strukturelt misforhold: «Så vi har ikke et slikt kundenummer, så det er bare en skrivefeil, ikke sant?». Erik fra Stena Recycling AS koblet disse feltene til organisatoriske krav og forklarte at: «Jeg måtte legge inn deklarasjonsnummer», og dette er et nødvendig felt for farlig avfall som ikke finnes i applikasjonen nå, og jeg håper at det finnes dersom jeg skal bruke applikasjonen.
Her konkluderer vi med at disse kvalitative resultatene viser grunnleggende feil som de kvantitative spørreundersøkelsene ikke klarte å avdekke. Selv om spørsmål fem, som gjelder metadata, fikk høye skårer (98,2\% og deretter 91,8\%), viste intervjuene noe annet. De viste at høy «opplevd kvalitet» kan skjule strukturelle svakheter i utformingen av feltene og deres egnethet for den praktiske konteksten.





\textbf{Teorikobling.} Den teoretiske forbindelsen: Denne drøftingen går utover bare å registrere vanlige feil og knytter dem til metoder for feilforebygging og konsistens \cite{nielsen1994}, samt til deres direkte innvirkning på dimensjonene ved datakvalitet når det gjelder fullstendighet, entydighet og korrekt kobling \cite{wang1996}. Generelt viser resultatene at feiltypene ikke oppsto tilfeldig, men følger påfølgende akkumuleringer som samsvarer med forskningens teoretiske modell; problematikken begynner med uklarhet i systemets tilstand → manglende samsvar med brukerens mentale modell → noe som fører brukeren til en feil handling → og dette ender til slutt i inkonsistens i dataene \cite{nielsen1994, norman2013,wang1996}.
Intervjuene viser at Geiger. \cite{geiger2021} sa at «menneskelig klassifisering er den øvre grensen for kvalitet», og dette er ikke begrenset til selve klassifiseringsprosessen. Intervjuene viste at dette prinsippet også gjelder kategoriene som systemet opprinnelig tilbyr brukeren. Dette bekrefter Heck observasjon \cite{heck2025}: Datakvalitet styres ikke i slutten av prosessen, men i hvert steg fra begynnelsen.

\textbf{Bidrag til problemstilling.} Poenget her ligger i hvor tydelig det tredje forskningsspørsmålet viser hvordan mekanismen kommer til uttrykk i problemer i brukergrensesnittet gjennom svakheter i brukergrensesnittet (UI) og mangler ved systemets tilstand, som fører til konkrete inkonsistenser i datastrukturen. Denne sammenhengen mellom observerbare problemer i brukergrensesnittet og deres innvirkning på datakvaliteten gjenspeiles empirisk gjennom fritekstkommentarene, og vil bli ytterligere tydeliggjort når analysen av tenke høyt blir tilgjengelig. Intervjuene avdekket et dypere problem: problemet er at klassifiseringsmodellen i appen og de tilknyttede beskrivende dataene ikke gjenspeiler hvordan arbeiderne faktisk håndterer avfallet. Dette gjør de innsamlede registreringene mindre nyttige senere for trening av kunstig intelligens, uavhengig av nøyaktigheten i selve inndataene.

\subsection{RQ4:- Hvilke konkrete forbedringer i design og implementasjon kan redusere feil og styrke datakvaliteten i registreringene?}

\textbf{Funnkilder.} Kvantitative resultater fra begge brukerundersøkelser, fritekstkommentarer, og dybdeintervjuer med Franser Stena, Erik Stena, Benjamin, Werner Reiersen Stena og NG Roger Aas.

\textbf{Funn – teoretisk bakgrunn: } 
Basert på planleggingsfasen var det forventet at de forventede forbedringene skulle formes på tre nivåer for å sikre systemets effektivitet og kvaliteten på dataene.

\textbf{Brukergrensesnitt:} Denne situasjonen fokuserte på å forbedre og utvikle statusindikatorene for å styrke ``systemets gjennomsiktighet'', gjennom tydeligheten i den aktive registreringsprosessen via klar synlighet, eksplisitt bekreftelse ved innsending, og muligheten til å se den nye registreringen etter at systemet er tilbakestilt på nytt.

\textbf{Målinger av arbeidsflytens forløp:} Det ble arbeidet med å forbedre og utvikle sekvensperspektivet, gjennom validering av data før innsending (Data Validation), og målet er å redusere den nødvendige mentale innsatsen.

\textbf{Gjennomføringsmålinger:} Det ble lagt vekt på datakontinuitet, inkludert det som omfatter lagring av utkast, og å ta i bruk en streng logikk for tilbakestilling (Reset Logic), samt tydelig håndtering av de to tilstandene ``før/etter innsending''.

\textbf{Forventet diskusjonsstruktur.}
\begin{enumerate}
  \item \textbf{UI-tiltak}: statusindikatorer, tydelighet rundt «aktiv registrering», eksplisitt bekreftelse ved innsending, og synlighet av «ny registrering» etter tilbakestilling.
  \item \textbf{Arbeidsflyt‑målinger}: rekkefølge, validering før innsending og reduksjon av kognitiv belastning
  \item \textbf{Implementasjons‑målinger:}: sterk datakontinuitet (inkludert utkastlagring), tydelig logikk for tilbakestilling og klar håndtering av «før/etter innsending»-tilstander.
\end{enumerate}

\textbf{Funn - empiriske resultater}

Resultatene viste tre tiltak med tydelige effekter i forbedringen av resultatindikatorene:

\begin{enumerate}
    \item Tekstetiketter under symbolene ble aktivert (+52,0 prosentpoeng i Q8), og dette tiltaket var det mest virkningsfulle når det gjaldt brukervennlighet og datakvalitet. Derfor regnes dette tiltaket som en direkte anvendelse av Nielsens heuristiske teori, som handler om prioriteringen av «gjenkjenning av elementer fremfor gjenhenting fra hukommelsen» \cite{nielsen1994}. I en oppgaveorientert applikasjon, der registreringen skjer i et arbeidsmiljø preget av press og begrensede oppmerksomhetsressurser \cite{muller2024, weichbroth2024}, reduserer denne visuelle forbedringen risikoen for feil som oppstår når brukeren trykker på feil knapp eller avbryter prosessen fordi vedkommende ikke finner sendeknappen.

    \item Den tydelige bekreftelsen av innsendingen (+22,7 prosentpoeng i Q7) bidro til å bryte den betingede sammenhengen mellom uklarhet i systemets tilstand og forekomsten av dobbel innsending. Den tydelige og synlige forbedringen etter innsending gjenspeiler også det Norman \cite{norman2013} omtaler som tilbakemelding, ved at systemet bekrefter at handlingen er fullført og at tilstandsovergangen er gjennomført. Dette bekrefter for oss at dette er et direkte mål på datakvalitet ved å eliminere tilstanden av usikkerhet som driver brukeren til å sende på nytt, og dermed bryter med kriteriet om entydighet som Wang og Strong har fastsatt \cite{wang1996}.

    \item Etter fullføringen av innsendingen sikrer den korrekte tilbakestillingslogikken (+17,2 prosentpoeng i Q3) at hvert bilde kobles korrekt til sin egen post, og ikke kopieres til den neste prosessen. Denne implementeringsdetaljen er direkte knyttet til korrekt kobling \cite{wang1996}: Dersom det ikke finnes et korrekt tilbakestillingssystem, vil applikasjonen produsere poster der gamle bilder kobles til nye metadata. Dette resultatet bekrefter at den tekniske håndteringen av tilstandsovergangen «sendt ← ny post» er viktigere for datakvaliteten enn selve utformingen av brukergrensesnittet.
\end{enumerate}

De kvalitative kommentarene i andre runde identifiserte prioriteringer for videre fremtidige forbedringer.

Flyttingen av knappen for hovedsiden var det hyppigst gjentatte temaet i kommentarene (6 av 17 kommentarer), og dette må tas i betraktning i universell utforming. Utformingen av kategorivelgeren må vurderes på nytt slik at den ikke skjuler innholdet bak tastaturet, for å sikre valg av kategori uten forstyrrelser som følge av bevegelse. Til slutt må påliteligheten i datakontinuiteten under navigasjon mellom skjermene verifiseres, i samsvar med anvendelsen av AsyncStorage-protokollen som er beskrevet i avsnitt 3.4.

\newline

\textbf{Funn fra dybdeintervjuene: } 
Intervjuene som ble gjennomført, bidro til å utvide listen over forbedringer betydelig, der disse forslagene ble klassifisert ut fra deres forventede påvirkning på datakvalitet og brukervennlighet:

\textbf{Forbedringer med stor effekt (viktige for å sikre datakvalitet):}

Franser fra Stena, Benjamin, Werner Reiersen fra Stena og Roger Aass (fra NG) foreslo å innføre egne klassifiseringer for hvert selskap separat, eller å tillate «hovedbrukeren» (Key User) å definere disse klassifiseringene. Gjennom dette forslaget løses problemet med klassifiseringene som er identifisert i det tredje forskningsspørsmålet, fordi hver organisasjon har sine egne forhold og operative prosedyrer. Den forventede effekten her vil også gjenspeiles i kvaliteten på klassifiseringene, økt nøyaktighet i kunstig intelligens-algoritmene og reduksjon av tilfeller med feil valg.

Roger Aass og Werner Reiersen fra Stena foreslo også å legge til manglende kategorier i listen, som «restavfall». Dette tiltaket regnes som en grunnleggende nødvendighet for å bruke applikasjonen i de daglige aktivitetene til selskaper som NG og Stena, der restavfall utgjør en av de viktigste kategoriene hos dem.

Protokoll for håndtering av farlig avfall: Erik fra Stena snakket om å aktivere et eget løp for farlig avfall, og påpekte at denne typen krever en egen organisatorisk plan, som «deklarasjonsnummer, avfallsstoffnummer og begrunnelse», som er informasjonsfelt som ikke dekkes av den nåværende metadata-strukturen i applikasjonen.

Erik og Werner Reiersen fra Stena, samt Roger Aass, var også enige om betydningen av å integrere applikasjonen gjennom et programmeringsgrensesnitt (API) med systemer for virksomhetsressursplanlegging (ERP), vekter, avfallsdata og det europeiske varenummersystemet (EAN). Dette forslaget regnes som svært komplekst fra et teknisk perspektiv, men har størst effekt. Roger Aass snakket om et europeisk varenummerregister (EAN) og produktdata, mens Erik fra Stena på den andre siden snakket om et programmeringsgrensesnitt (REST API) tilpasset ERP-systemer, og Werner Reiersen fra Stena foreslo direkte kobling til avfallsdatabaser. Dette tiltaket flytter applikasjonen fra å være et selvstendig registreringsverktøy til å bli en organisk del av et integrert system for datainnsamling.

Mekanisme for midlertidig lagring (arbeid uten tilkobling): Benjamin og Werner Reiersen fra Stena foreslo å arbeide med å etablere en funksjon for midlertidig lagring av data når internettforbindelsen brytes (Offline Caching), slik at de sendes automatisk når nettverket gjenopprettes. Werner Reiersen fra Stena uttrykte denne operative forventningen tydelig ved å si: «Jeg forventer at dataene lagres midlertidig», og dette er en svært viktig funksjon ved arbeid med avfall som kan ha blitt flyttet, og da kan vi ikke registrere informasjonen på nytt, noe som er svært viktig.

Utvikling av hjelpetekster (Help Texts): Benjamin og Roger Aass foreslo å tilby tydeligere forklarende tekster for både «ordre-ID» og «kunde-ID», slik at uklarheten rundt metadatafeltene, som tidligere ble diagnostisert i det tredje forskningsspørsmålet (RQ3), fjernes.

\textbf{Forbedringer med middels effekt (brukervennlighet og kontroll):}

Rekonstruksjon av det visuelle grensesnittet: Erik, Werner Reiersen Stena og Benjamin snakket om å flytte knappen for hovedsiden til midten av skjermen, og dette forslaget samsvarer med de kvantitative resultatene fra spørreundersøkelsen. Benjamin påpekte også at profilknappen som for øyeblikket er plassert i midten, er den minst brukte i praksis, og at vi kan erstatte den med noe viktigere på grunn av at profilikonet allerede finnes øverst på siden, mens Werner Reiersen fra Stena kommenterte: «Mange applikasjoner plasserer hovedsiden i midten», og Erik Stena oppsummerte sitt syn enkelt: «Jeg foretrekker å plassere den viktigste knappen i midten».

Å gi fleksibilitet i kontroll og brukerfrihet: Werner Reiersen fra Stena la frem et forslag som gjør det mulig å trekke tilbake en bestilling eller endre vinduet før fakturaen utstedes, noe som gir brukeren et større område for kontroll og frihet \cite{nielsen1994}, og reduserer bekymringer knyttet til generering av feil fakturaer.

Erik Stena, Benjamin og Roger Aass foreslo å vise de tatt bildene på siden for historikk og detaljer. Benjamin begrunnet dette ved å si: «Du kan gå tilbake til det bildet senere», og dette er et tiltak som sikrer en «korrekt kobling» \cite{wang1996} som brukeren kan kontrollere.

Forenkling av oppstartsprosedyrer: Werner Reiersen og Erik fra Stena foreslo å kreve aktivering av alternativet «glemt passord» for å overvinne innledende innloggingsproblemer.

\textbf{Forbedringer med lav til middels effekt (dekning av tilpassede brukstilfeller):}

Denne kategorien omfattet en gruppe supplerende forslag som tjener bestemte brukergrupper, og består av: å legge til søke- og filtreringsfunksjon i nettleserhistorikken (Benjamin), og støtte for stor tekst for å styrke tilgjengeligheten (Werner Reiersen fra Stena). Disse endringene bidrar samlet sett til å forbedre plattformens generelle effektivitet og tilpasse den til bestemte arbeidsforhold.


\textbf{Teorikobling.} Forbedringene som ble gjort i applikasjonen, kan begrunnes med henvisning til det heuristiske prinsippet, eller til den delen av forskningsmodellen som omhandler dette, nemlig å styrke transparensen i systemets tilstand \cite{nielsen1994} og redusere gapet mellom systemet og brukerens mentale modell \cite{norman2013}. Disse resultatene viser at datakvalitet i systemer som opprettes av brukerne ikke bare er et algoritmisk problem, men begynner i det øyeblikket brukeren samhandler med innsamlingsgrensesnittet, slik både Geiger mfl. \cite{geiger2021} og Heik \cite{heck2025} har vist.
Offline-buffer adresserer error prevention \cite{nielsen1994} og styrker completeness \cite{wang1996}. ERP-integrasjon reduserer manuell registrering og dermed hele kjeden av menneskeskapte feil som dokumenteres hos Geiger et al. \cite{geiger2021}. Tilbaketrekking før fakturering implementerer user control and freedom \cite{nielsen1994}.


\textbf{Bidrag til problemstilling.} Det fjerde forskningsspørsmålet (RQ4) ligger i det anvendte bidraget i denne drøftingen: hvilke endringer kan bidra til å øke datakonsistensen gjennom å forbedre brukeropplevelsen og håndtere systemets tilstand på en bedre måte? De tre dokumenterte tiltakene utgjør også et bevis på erfaringsbaserte designanbefalinger, som kan anvendes direkte i utviklingen av StellAI.
\newline Resultatene fra disse intervjuene bidrar til designutvikling som inkluderer strukturelle og grunnleggende forbedringer utover bare visuelle modifikasjoner. Det har dukket opp et kritisk behov for å ta i bruk kontekstuelle kategorier, ERP-integrasjon og aktivering av offline caching-mekanismer. Disse forbedringene er viktige pilarer for å sikre applikasjonens effektivitet i virkelige arbeidsmiljøer. Dette inkluderer å tilby data av høy kvalitet og pålitelighet som er tilstrekkelig for å trene AI-modeller og transformere systemet fra et enkelt datainnsamlingsverktøy til et omfattende teknologisk system som støtter beslutningstaking.


%\section{Konklusjonen: fra forskningsspørsmålene til behandling av problemstillingen}
%Etter gjennomføringen av brukertestene bør drøftingen avsluttes med en oppsummering som tydelig besvarer forskningsspørsmålet ved å integrere spørsmål RQ1 til RQ4:

%\begin{itemize}
% \item RQ1 identifiserer utfordringene brukerne møtte og uklarheten knyttet til systemtilstanden.
%  \item RQ2 forklarer mekanismen (mentale modeller og manglende samsvar) som fører til feilaktige handlinger.
%  \item RQ3 identifiserer feilene og datainkonsistensene som oppstår som følge av denne mekanismen.
%  \item RQ4 oppfordrer til endringer i design og implementering for å adressere de identifiserte årsakene, noe som bidrar til å forbedre datakonsistensen.
%\end{itemize}
%Oppsummeringen vil følge den kausale sekvensen i forskningsmodellen, og gjøre sammenhengen mellom brukergrensesnitt/brukeropplevelse, systemtilstand og datakonsistens tydelig og etterprøvbar /%\cite{nielsen1994,norman2013,wang1996}.